\section{The generic principal-character estimate}
\label{sec:generic-character-complete-assembly}

We now instantiate the generic estimate
\eqref{eq:sr-generic-count}.
The auxiliary imaginary quadratic field is always
\(K=\mathbb Q(\sqrt{-3})\).
No restriction on splitting of \(\ell\) in \(K\)
is imposed.

\subsection{Classicality at every dominant weight, including \(a=1\)}
\begin{srlemma}\label{lem:sr-good-point-classical-all-dominant}
The \(\mathrm{GoodPt}_n\) conclusion of
Proposition~\ref{prop:actual-common-two-three-endpoint}
gives the actual ordinary automorphic point at every
dominant algebraic specialization of its weight ring.
This includes rank three at exponents
\((a-1,-1,-a-1)\) for every \(a\ge1\).
\end{srlemma}
\begin{proof}
Let \(\mathcal M_n\) be the finite free ordinary weight module
and \(\mathbb T_n\subseteq\operatorname{End}_{\Lambda_n}\mathcal M_n\)
its faithful Hecke image.
It is also finite as a \(\mathbb T_n\)-module.
Its support over \(\mathbb T_n\) is all of
\(\operatorname{Spec}\mathbb T_n\), because its
annihilator is zero.
A field-valued Hecke point \(x\), supplied by
\(\mathrm{GoodPt}_n\), therefore has
\[
                       \mathcal M_n\otimes_{\mathbb T_n,x}E\ne0 .
 \tag{GC1}
\]
Indeed localize at \(\ker x\), apply Nakayama, and
then extend its residue field to \(E\).
In particular its character occurs in the corresponding
specialization \(\mathcal M_n\otimes_{\Lambda_n}E\).

Proposition~\ref{prop:classical-ordinary-control}
applies at every dominant algebraic weight.
It requires neither strict dominance of the algebraic
highest weight nor a lower bound on \(\ell\) in terms
of that weight.
Here is the boundary check in the exact normalization.
For ordered Galois exponents \(h_i\), equations
(CA3)--(CA4) give algebraic weight
\(\lambda_i=h_i+i\), up to the harmless common
scalar weight and the fixed choice of dual Borel.
For rank three this is
\[
           \lambda=(a-1,0,1-a),\qquad
           \lambda_i-\lambda_{i+1}=a-1\ge0 .
 \tag{GC2}
\]
The controlling positive operator has ordinary slope
zero, strictly below
\(\lambda_i-\lambda_{i+1}+1=a\).
Equivalently, in the integral proof (CC4)--(CC10),
every nonextremal weight of the Schur module differs
from the extremal weight by a nonzero sum of positive
roots. A strictly positive diagonal element contracts
it by at least one coefficient uniformizer.
Zero simple-root differences merely remove some
of these weights; they do not produce additional
uncontracted coordinates.
At \(a=1\), \(\lambda=(0,0,0)\) and the algebraic
representation is one-dimensional, so those
other coordinates are absent and the control
map is immediate.
For rank two the adjacent difference is
\(2a-1\ge1\), unchanged by the scalar link character.
These checks apply to the Gaussian rank-three
specialization as well.

Thus (GC1) is a nonzero quotient of the actual
classical ordinary space with the same
\(\tau_0,\theta^{-1}\) and normalized \(U_j\).
Choose a classical automorphic constituent carrying
this Hecke character.
The classical Galois realization of
Section~\ref{sec:common-ordinary-arithmetic-action}
has its determinant.
The given point and that semisimple representation
are isomorphic by Chebotarev and Brauer--Nesbitt.
If the point on the selected CM group is irreducible,
the isobaric base change of this constituent has
only one cuspidal block of full rank:
any other discrete-block decomposition gives a
reducible semisimple Galois representation.
This proves the asserted cuspidal realization.
The use of very regular weights to separate the
Hecke algebra in (CA6) is therefore distinct from,
and does not restrict, this all-dominant classicality.
\end{proof}

For type I, the actual ordinary flag and the
nonresonance condition prove irreducibility on
every open subgroup, as in
Section~\ref{sec:type-one-rigidity}.
The actual seed tower therefore preserves
irreducibility of these points.
For the rank-two outer block, its two residual
characters restrict distinctly to \(G_L\):
their ratio is \(\bar\Psi^2\ne1\) on a selected
ordinary local group, by (SP16).
The finite extension \(L/K\) is abelian.
Lemma~\ref{lem:sr-residual-distinct-abelian-restriction},
applied in rank two over \(K\), consequently
preserves the irreducibility of the outer block
on this actual tower.
Thus the irreducibility premise in (GC1)
has been verified in both runs.

Remove the algebraic link character by
Lemma~\ref{lem:sr-link-character-algebraic}.
Exact cyclic descent through the CM tower,
as in Lemma~\ref{lem:sr-cyclic-exact-descent},
then gives the fixed-field rank-three endpoint
of Proposition~\ref{prop:sr-generic-three-exact-endpoint}.
The local result
Corollary~\ref{cor:sr-generic-original-level}
gives \(\mathrm{Th.Unl}\) and the original
\(K_3(N)\)-fixed vector, including at \(a=1\)
and at \(\ell\mid N\).
The final representation has weight
\[
                         (a-1,-1,-a-1),                    \tag{GC3}
\]
not the arbitrarily large weight used to construct
the initial Eisenstein seed.

The analogous rank-two endpoint is rational.
We fix its normalization before descending.
The outer representation \(W\) is defined on
\(G_{\mathbb Q}\), has determinant \(\eta_K\),
and its restriction to \(G_K\) is irreducible.
After removing the algebraic CM link, restore
the algebraic and half-norm normalization of
the rank-two correspondence before applying
quadratic cyclic descent.
Equivalently, apply descent to the system
normalized to attach to \(W(-a)=W\epsilon_\ell^{-a}\).
Its unramified geometric-Frobenius roots are
\(p^a\) times the roots of \(W\);
this prescription fixes the scalar normalization
at every unramified place.
We make no assertion that the unshifted
unitary parameter of \(W\) is itself an
unshifted regular-algebraic automorphic representation.

The normalized system over \(K\) is invariant
under \(\operatorname{Gal}(K/\mathbb Q)\):
the Galois representation extends to
\(G_{\mathbb Q}\), and strong multiplicity
one identifies its conjugate system.
Quadratic cyclic descent is cuspidal, since
the original restriction is irreducible.
Its real parameter is the holomorphic
discrete-series parameter of weight \(2a+1\),
with the scalar fixed by the preceding
unramified normalization.
Newform theory therefore gives a primitive
holomorphic eigenform \(f\) of that weight.
The exact newform attachment used in
Section~\ref{sec:kato-rigidity} is
\[
 \det\!\left(X-W(-a)(\operatorname{Frob}^{\rm geom}_p)\right)
       =X^2-a_p(f)X+\omega_f(p)p^{2a}
       \qquad(p\nmid 3N\ell).
 \tag{GC3a}
\]
Thus the Tate twist here is the inverse
cyclotomic twist; \(\epsilon_\ell
(\operatorname{Frob}^{\rm geom}_p)=p^{-1}\).
The half-norm in algebraic local Langlands
has already been restored when forming this
newform system and is not discarded during
the CM descent.

Deligne's representation of \(f\) restricts
to \(W(-a)|_{G_K}\).
Schur's lemma makes its possible discrepancy
from \(W(-a)\) a character of
\(\operatorname{Gal}(K/\mathbb Q)\).
Twisting the descended form by this finite
quadratic character corrects the discrepancy
exactly. This preserves the weight and
cuspidality. The resulting primitive newform
satisfies (GC3a) with \(W(-a)\) itself.
In particular \(\omega_f=\eta_K\);
the finite nebentype and the factor \(p^{2a}\)
also fix the scalar normalization independently
of the angular description at infinity.

Full local comparison and the inherited local
flags give the rational witness
\(\mathrm{WitAQ}\).
Its finite inertia characters have conductor
bounded by the original \(3N\), with \(N=0\)
for their Weil--Deligne monodromy.
At \(\ell\) this also follows by taking the
outer subrepresentation of the crystalline
representation on the subgroup in (LC13);
the exponent difference is \(2a\ge2\).
At three the possible extra character is
exactly \(\eta_K\).
The induced principal-series vectors therefore
give the principal congruence level \(3N\).
These are all the clauses of
\(\mathrm{Adm2A}\), not merely modularity at
an unspecified level.

These arguments prove the actual
\srcid{AutomorphyUnramifiedCoefficientsAlt}
output used in \(\mathrm{FOneCharU}\):
one common marking, finite weight-zero
algebra, type-I fixed-field automorphy,
and modularity of the rank-two extensions.
The characteristic-\(\ell\) finite-part assertion
is Proposition~\ref{prop:actual-common-two-three-endpoint}
and Proposition~\ref{prop:sr-host-finite-part-from-finend}.
Classicality and exact descent are the preceding
arguments. The finite integral representative
of a geometric point is obtained by the graded
torus scaling already proved in
Section~\ref{sec:spectrum-deformation}.

\subsection{Rigidity and the one-character inequality}
Fix \(\ell\notin\{2,3,7\}\), \(a\ge1\) with
\(\ell-1\nmid2a\), and an admitted host character
\(\psi\).
The good-data hypotheses of
Section~\ref{sec:single-character} hold by
Proposition~\ref{prop:sr-three-auxiliary}.
In particular, if \(\ell\mid N\), write
\[
 \psi=\psi^{(\ell)}\widetilde\omega_\ell^{\,b},
 \qquad b=0\ \hbox{or}\ (\ell-1)/2.
\]
Then
\(\bar\Psi^2|_{I_\ell}=\bar\epsilon_\ell^{\,2a}\ne1\).
This is exactly the residual separator used
above. Absorbing \(b\) in the auxiliary seed
weight changes no weight in (GC3).

There is also no inertness requirement here.
For split \(\ell\) the sign \(\eta_K\)
is already trivial on its decomposition group.
For inert \(\ell\ne3\) it is the unramified
quadratic sign. In either case the constructed
tower has split ordinary places above \(\ell\),
and \(H=\ker\eta_K\cap\ker\tau\) in (LC13)
is the appropriate open local subgroup.
The proof of local crystallinity and descent
of the ordered flag applies in both cases.

At a type-I mixed point,
Proposition~\ref{prop:sr-cyclotomic-restriction}
gives irreducibility on
\(G_{K(\mu_{\ell^\infty})}\).
The tangent module over \(\mathbb Q\) is
the descended polarized adjoint:
on \(G_K\) it is ordinary conjugation and
on complex conjugation it is
\[
                         X\longmapsto-J_3X^tJ_3^{-1}.
\]
The symmetric pairing and multiplier
\(\epsilon_\ell^{-2}\) are precisely those
of (GE3). Hence
\href{https://arxiv.org/pdf/2207.04925}
{Thorne, Theorem 6.1} gives its
Bloch--Kato Selmer vanishing over \(\mathbb Q\).
This is the descended adjoint in that theorem;
no full-adjoint vanishing over \(K\) is substituted.

The actual tangent cocycles satisfy those
local conditions.
Away from \(\ell\), the represented killed-inertia
condition makes their restrictions inertia
coboundaries. At \(\ell\), (LC8)--(LC14)
applies also to dual-number points of the
actual represented local condition.
Their crystalline extensions therefore give
the finite local condition.
Restriction to the finite local extension is
legitimate; averaging the corresponding local
primitive over the finite quotient gives the
original finite local class.
Thus the published theorem acts on these
actual tangents.

For a type-II point the tangent decomposition
proved in Section~\ref{sec:kato-rigidity} is
\[
 \operatorname{ad}r\simeq
    \operatorname{ad}^0W\oplus W\oplus W\eta_K
                         \oplus\eta_K\oplus\eta_K .
 \tag{GC4}
\]
The actual rational newform just constructed
has weight \(2a+1\) and is attached with the
exact twist \(W\epsilon_\ell^{-a}\).
Kato's finite Selmer theorem at the critical
integer \(a\), as applied in that section,
gives the finite-class vanishing for \(W\)
and \(W\eta_K\); the inequalities are
\(1\le a\le2a\) and \(2a\ne2a+1\).

For the adjoint summand one may verify
Newton--Thorne's hypothesis directly.
A nontrivial quadratic self-twist of \(W\)
preserves the unique higher-Hodge-weight line
of its ordinary local flag. It acts trivially
on its one-dimensional quotient, so the
twisting character is trivial on the entire
local decomposition group at \(\ell\).
A CM quadratic field contained in
\(\mathbb Q(\mu_{\ell^\infty})\), for odd
\(\ell\), is ramified at \(\ell\).
Its self-twist would contradict this
local triviality. Therefore the CM exception
in Newton--Thorne, Theorem 5.4, does not occur;
the non-CM case is already covered by that
theorem. It follows that
\(H_f^1(\mathbb Q,\operatorname{ad}^0W)=0\).
Finally \(\eta_K\) becomes trivial over \(K\);
the finite local condition for the trivial
character is unramified everywhere, and
global class field theory factors it through
the finite class group. Thus its vector
Selmer group is zero. Equation (GC4) proves
type-II tangent vanishing as well.

At the reducible base point the same
nonpositive-weight local finite-class
calculation gives the prescribed triangular
tangent statements.
For exponent zero use the finite class group.
For negative exponents, the local finite
condition is zero at \(\ell\); the global
strict vector group is zero by the finiteness
comparison with \(K_{2b}\) already established
in Section~\ref{sec:sr-character-finiteness}.
Only finiteness, not a polynomial order
estimate, is used at this step.
The statement applies to the finite abelian
character fields involved here, including
the fixed quadratic sign.

All algebraic inputs of the one-character
proof are now instantiated: the marked
represented condition and its finite tangent
and obstruction spaces, normalized
Poitou--Tate duality and global Euler
characteristic, the proved uniform
\(H^2\) bound, the finite weight-zero
algebra, the actual mixed-point
realizations, and the just-checked
tangent vanishings.
The lower cotangent comparison and the
graded intersection calculation of
Section~\ref{sec:single-character}
therefore give the following actual assertion.
For every \(\sigma\in G_K\), if
\[
  \operatorname{tr}\rho(\sigma)\ne
       \Psi(\sigma)+1+\Psi(\sigma)^{-1}
\]
for every member of the two admissible
classes, there is a finite set \(X_\psi\)
of pairwise nonisomorphic admissible
representations such that \(V'_{\ell,N,a}(\psi)\)
is finite and
\[
 v_{\ell,a}(\psi)\le
 2\sum_{\rho\in X_\psi}
 v_\ell\!\left(\operatorname{tr}\rho(\sigma)
                   -\Psi(\sigma)-1-\Psi(\sigma)^{-1}\right).
 \tag{GC5}
\]
The full decomposition conditions at
\(\ell N\) and all finite inertia conditions
in this assertion are exactly the defining
conditions of \(V'\).

The coefficient normalization is unchanged.
First realize all prime-to-\(\ell\) character
values in an unramified coefficient DVR.
For two such realizations and an arbitrary
finite splitting DVR take their finite
compositum. Cohomology, restrictions and
their strict kernels commute with this finite
free coefficient extension.
The module length and
\(v_{\mathcal O}(\ell)\) both multiply
by the relative ramification index.
Thus (GC5), including its finiteness
assertion, holds for the coefficient DVR
in the question.
This is the exact normalized quantity
\(\operatorname{length}_{\mathcal O}(V')/
  v_{\mathcal O}(\ell)\).

\subsection{A coefficient-independent finite spectrum}
We make the fixed constants in the trace
integer independent of the large support
tower. Only the original quadratic field
and original arithmetic weights occur here.

For rank three use the fixed field
\(K=\mathbb Q(\sqrt{-3})\), level \(N\)
and weight (GC3).
The coefficient construction in
Section~\ref{sec:spectrum-deformation}
has dimension at most \(C(a+1)^{18}\).
The proved arithmetic cohomology bound
Proposition~\ref{prop:sr-level-cohomology}
and the injection of cuspidal cohomology
therefore bound its pairwise nonisomorphic
systems by
\[
                           C_3[N(a+1)]^{18}.               \tag{GC6}
\]
The constants precede \(N,a,\ell,\psi\).
The Lie-algebra cohomology criterion is the
regular-algebraic criterion of Clozel,
Theorem 3.13 and Lemma 3.14; cuspidal
cohomology injects into ordinary cohomology
by Borel,
\href{https://doi.org/10.1215/S0012-7094-83-05028-7}
{\emph{Regularization theorems in Lie algebra
cohomology. Applications}}.
Their hypotheses are exactly the regular
algebraic cuspidal representations and
the dual coefficient system constructed
there.
In particular these theorems are applied
to the fixed field \(K\), not the growing
totally real tower.

For the rational rank-two systems one can
retain the fixed quadratic catalogue of
Section~\ref{sec:common-trace-integer}.
Take
\[
 K_i=\mathbb Q(\sqrt{-d_i}),\qquad
              d_i=3,7,11,19,23,31,
\]
and omit the field of discriminant \(-7\)
when the test prime is seven.
At most four quadratic self-twists of a
rank-two representation exist, by the
endomorphism-space argument of
Lemma~\ref{lem:sr-quadratic-self-twists}.
Thus at least one of the five remaining
quadratic base changes of each rational
cuspidal system is cuspidal.
The same statement follows from cyclic
cuspidality and its Galois realization.
A fixed cuspidal base-change class has
at most two rational preimages, the two
twists by the quotient character.
All these assertions concern a fixed
finite list of fields.

At these fixed CM fields the inherited
principal congruence level is \(3N\).
This can also be checked character by
character: for a local quadratic extension
of ramification index \(e\), the norm of
\(1+\mathfrak p_L^{ec}\) lies in
\(1+\mathfrak p_{\mathbb Q_p}^{c}\).
The finite inertia characters that were
trivial at level \(p^c\) therefore remain
trivial at the scalar level \(p^c\) over
the quadratic field. The monodromy is
zero and the principal-series vector
has that congruence invariance.
The infinitesimal weights are the fixed
rank-two weights \((a,-a)\), in the
normalization of the rational trace
construction.
The same coefficient/cohomology argument
has exponent \(8=2\cdot2^2\) in rank two.
Maximize its constants over the five
fixed fields, and include the factor two
from the fibers. The rational systems
are consequently bounded by
\[
                            C_2[N(a+1)]^8.                \tag{GC7}
\]
It suffices to take the systems which
arise from the primitive holomorphic
newforms just obtained. This is a
coefficient-independent definition of
the label set; the larger CM
cohomological spaces give its bound.
No requirement that a catalogue field
be inert at the varying \(\ell\) is
used in (GC7).

The finite complex label sets in
(GC6)--(GC7) are stable under every
automorphism of \(\mathbb C\).
For rank three this is Clozel's
rationality/conjugation theorem,
Theorem 3.13: the finite part is defined
over a number field and its conjugates
are again regular algebraic cuspidal,
with the conjugate algebraic weight.
Our weights are parallel, so they are
unchanged; level and polarization are
preserved.
For rank two this is the rational
structure on the holomorphic newform
space, equivalently the same theorem
with the conventional half-norm shift.
All Satake roots and traces below use
the integral normalization, so that
shift introduces no ambiguity.
The disjoint union of the two label
sets has size at most
\[
                     C_*[N(a+1)]^{C_*'}                   \tag{GC8}
\]
for constants fixed before all hosts
and coefficient primes.
Equality of global labels makes the
attached semisimple representations
isomorphic, by Chebotarev and
Brauer--Nesbitt. Labels may have equal
numerical traces at the test place;
they remain separate labels in the product.

\paragraph{Scaled integral roots at the test prime.}
There is no additional integral Satake
hypothesis left in this application.
Put \(p=7\), choose \(w\mid7\) in \(K\),
and set \(q_0=Nw\), \(Q=q_0^a\).
Here \(q_0=7\), since seven splits in
\(\mathbb Q(\sqrt{-3})\); the following
argument also works with residue degree
two.

For each rank-three label, let
\(\alpha_1,\alpha_2,\alpha_3\) be the
roots for the original untwisted
representation at geometric Frobenius.
They are algebraic by rationality.
At a rational prime \(r\ne7\), the
polarized Galois realization at \(r\)
has compact image and hence a stable
lattice over a finite coefficient field.
Its Frobenius roots are integral at
every place above \(r\).
The scalar inverse-cyclotomic twist
has unit value there, so the same
is true of the \(\alpha_i\).
For a place above seven use instead
the crystalline realization at seven,
which exists since the local
automorphic component is spherical.
Its untwisted exponents are
\((a,0,-a)\).
The Newton--Hodge inequality gives
\[
                       v_7(\alpha_i)\ge-a\,v_7(q_0).
\]
Indeed the least Newton slope has a
Frobenius-stable isoclinic submodule;
its induced Hodge slopes are bounded
below by \(-a\), and weak admissibility
gives the displayed inequality.
Full local compatibility identifies
these crystalline Frobenius roots
with the same algebraic Satake roots.
Thus \(Q\alpha_i\) is integral at
every finite place, hence is an
algebraic integer.

The polarized realization and
crystalline compatibility just used
have no residual image assumption:
they are the classical realization
and unramified local assertions of
Barnet-Lamb--Gee--Geraghty--Taylor,
Theorem 2.1.1, together with its
\(\ell=p\) compatibility in odd rank.
For rational rank two the same
integrality follows directly from
Deligne's newform representation
(or from its integral Hecke
polynomial); after normalization
its roots \(\beta_i\) satisfy
\(p^a\beta_i\) integral.
Over \(K\) the roots are their
residue-degree powers.
Consequently every trace label \(t_s\)
satisfies
\[
                         Qt_s\text{ algebraic integral}. \tag{GC9}
\]

Unitarize the cusp representation.
The strict unramified unitary Satake
bound puts each normalized root
between \(q_0^{-1/2}\) and \(q_0^{1/2}\)
in absolute value, and their product
has absolute value one.
The elementary three-root estimate
of Section~\ref{sec:common-trace-integer}
then gives
\[
                            |t_s|\le\sqrt{q_0}+2 .            \tag{GC10}
\]
For rank two append the constant
root \(1\); the same bound follows.
This is the Jacquet--Shalika
unramified bound in the exact
unitary normalization used there.
Equations (GC8)--(GC10) verify all
three hypotheses of the common
trace-product construction.

\subsection{The simultaneous prime sum}
\begin{srtheorem}\label{thm:sr-generic-character-estimate}
There is a natural number \(A\), fixed before
the principal host, \(t\), \(a\), the coefficient
DVRs, and the finite prime set, with the following
property.
For every sufficiently large principal host
\[
 E=\mathbb Q\bigl(
   \zeta_{2^{h+2}}+\zeta_{2^{h+2}}^{-1},
                  \sqrt{-q},\sqrt{-d}\bigr),\qquad
 N=2^{h+2}qd,
\]
with its stated principal congruence and size
conditions, every \(t\ge2\), every \(1\le a<t\),
and every finite set \(S\) of primes satisfying
\[
          \ell\notin\{2,3,7\},\qquad \ell-1\nmid2a,
\]
all the groups \(V'_{\ell,N,a}(\psi)\) in question
are finite, and
\[
 \sum_{\ell\in S}(\log\ell)
        \sum_{\substack{\psi\text{ a host character}\\
                        \psi(-1)=(-1)^a}}
 \frac{\operatorname{length}_{\mathcal O_\ell}
                    V'_{\ell,N,a}(\psi)}
                  {v_{\mathcal O_\ell}(\ell)}
                              \le (Nt)^A .                \tag{GC11}
\]
The splitting coefficient DVR \(\mathcal O_\ell\)
may be chosen independently at each \(\ell\).
\end{srtheorem}

\begin{proof}
Increase the fixed lower threshold once to
ensure the two principal odd primes exceed
seven.
For an orbit of characters with kernel \(J\)
in \((\mathbb Z/N)^\times\), use the common
label set just constructed and the roots
\(z_\psi=\psi(q_0)^{-1}\).
Its Eisenstein trace at inverse Frobenius is
\[
                 T_\psi=Q^{-1}z_\psi+1+Qz_\psi^{-1}.
\]
Since \(Q\ge q_0\ge7\),
\[
            |T_\psi|\ge Q-1-Q^{-1}>\sqrt{q_0}+2 .
\]
Thus (GC10) separates every admissible
trace, and (GC5) applies at this inverse
Frobenius. It proves finiteness before
any length is treated as a real number.

Equations (GC8)--(GC10) give the integer
\(Z_J\ne0\) of
Proposition~\ref{prop:sr-orbit-integer-implication}
with constants \(c,C\in\mathbb N\) fixed
only by the fixed spectrum:
\[
 \log|Z_J|\le
       c\,\varphi([(\mathbb Z/N)^\times:J])
                      [N(a+1)]^C\log7,\qquad
 \sum_{\ker\psi=J}v_{\ell,a}(\psi)
                            \le2\operatorname{ord}_\ell Z_J .
 \tag{GC12}
\]
Each integer is selected before \(\ell\) or
its coefficient DVR. The same integer controls
every allowed \(\ell\ne7\), because it is the
product of the full complex label set.
Good data hold for all members of a kernel
orbit, including the trivial orbit if admitted:
their orders are powers of two, primitive
powers preserve parity and the local quadratic
components, and the fixed field is \(K\).
There is no exclusion of the primes dividing
\(N\); their finite quadratic component was
already included in (GC2) and (LC13).

The kernel orbits partition \(\mathcal X_a\)
and, since the host Galois group is abelian
of order \(2^{h+2}\),
\[
       \sum_{J}\varphi([(\mathbb Z/N)^\times:J])
                    =|\mathcal X_a|
                    \le[E:\mathbb Q]\le N .
\]
Sum (GC12) over these orbits and over the
arbitrary finite set \(S\).
Prime factorization of the nonzero integers
gives
\[
 \begin{split}
 \sum_{\ell\in S}\log\ell\sum_{\psi\in\mathcal X_a}
                              v_{\ell,a}(\psi)
 &\le2\sum_J\sum_{\ell\in S}
                  (\log\ell)\operatorname{ord}_\ell Z_J\\
 &\le2\sum_J\log|Z_J|\\
 &\le2c\log7\,N[N(a+1)]^C .
 \end{split}
\]
There is no factor \(|S|\).
Choose \(k\in\mathbb N\) once with
\(2^k\ge\max(1,2c\log7)\), and take
\(A=C+1+k\).
As \(a<t\), \(N(a+1)\le Nt\), \(N\le Nt\),
and \(Nt\ge2\). The last bound is at most
\((Nt)^A\), proving (GC11).
All constant choices precede the quantified
host and coefficient data.
\end{proof}

\paragraph{Scope and producer closure.}
The numerical orbit argument was previously an
implication from a trace integer and a one-character
inequality.
The actual support endpoint, all-dominant
classicality, fixed-field descent, local crystalline
component and original level, polarized tangent
rigidity, fixed-field spectral count, and scaled
integrality have now supplied those inputs for
the generic principal packet.
The source assembly is
\(\mathrm{FOneCharU}\) followed by
\(\mathrm{FPropGKU}\), with precisely
\(\mathrm{Adm3U_s}\) and \(\mathrm{Adm2A}\).
The proof of (GC11) is confined to the generic
estimate. It asserts no exceptional/resonant
character estimate and no completion of the
separate geometric comparisons in the manuscript.
